\section{The actual Xi-jet unit, its complete Hermite remainder and the conductor pullback}
\label{sec:actual-unit-hermite}

The original arithmetic amplitude in PQO1--3 and OMV's original-domain
paragraph is the value of $2\xi/h$ at an actual simple quartet,
not an independently chosen phase \cite{Cumulative}.
The complete source is the preserved accepted782 cumulative LaTeX:
PQO1--3, source lines 1430--1481, and OMV's original-domain paragraph,
source lines 2750--2787.
The exact original period conjugation is PCL1--10
\cite{PCLOriginal}, with its
\href{https://github.com/KokunoYumeto/zeta-function-research-reader/blob/487253b5b01851da5b9b98af4e146fecf1fab902/workbenches/splitzero-tandem/continuations/20260919-native-conductor-coprimality/providers/PCL_COMPLETE.tex#L19}{pinned complete proof}.
RLA2--10 and RLA21--35 \cite{RLA008} give the geometric family and
its exact local branches, maps and original metrics; that complete
source accompanies this continuation.
The following calculation constructs the pullback from the actual
Xi values and derivatives and retains the nonzero-value residual.
Every matrix inverse used below is displayed and checked algebraically.
Hermite interpolation and the Chinese-remainder coordinates are
classical constructions; their full formulas and proofs here specify
the original programme objects and do not claim those constructions
as new general results.  The completed zeta function is Riemann's
$\xi$ in the original theta-integral convention PQO uses, with
$g=2\xi$ retained throughout.  The retained human-source account
of its reflection and completion is T. M. Apostol's DLMF
25.4.3--25.4.4 \cite{DLMFXi}; the exact theta identity used here
is the one explicitly invoked by PQO3, not a new zeta identity.

\subsection{Two squared nodes and exact division without a zero assumption}

Let $0<\delta<1/2$ and $\gamma>2$ be arbitrary real parameters in
the original domain.  Retain the four labels and coordinates
\[
\begin{gathered}
 w_+=\delta+i\gamma,\quad w_-=\delta-i\gamma,\quad
 (w_{\mathrm M},w_{\mathrm N},w_{\mathrm T},w_{\mathrm E})
                  =(w_+,w_-,-w_-,-w_+),\qquad \rho_a=1/2+w_a,\\
 \lambda_+=w_+^2,\quad\lambda_-=w_-^2,\quad
 d=\lambda_+-\lambda_-=4i\delta\gamma\ne0,\\
 \beta=2(\gamma^2-\delta^2)=-(\lambda_++\lambda_-),\quad
 \eta=(\delta^2+\gamma^2)^2=\lambda_+\lambda_-,\\
 H(y)=(y-\lambda_+)(y-\lambda_-)=y^2+\beta y+\eta,\quad
 h(s)=\prod_a(s-\rho_a)=H((s-1/2)^2).
\end{gathered}\tag{AUH1}
\]
The letter $d$ here denotes the nonzero squared-node difference, not
a cutoff or a conductor coefficient.  The symmetry and entire-function
identities from the original theta integral, recalled in PQO3, are
$g(1-s)=g(s)$ and $g(\bar s)=\overline{g(s)}$.
Consequently there is a unique entire function $F$ with real Taylor
coefficients such that
\[
 g(1/2+w)=F(w^2),\qquad g=2\xi.
 \tag{AUH2}
\]
For a proof, expand the entire left side at zero.  Its odd coefficients
vanish by the first symmetry and its even coefficients are real by the
second.  If those coefficients are $c_{2n}$, define
$F(y)=\sum_{n\ge0}c_{2n}y^n$.  Cauchy's coefficient bound on a
circle $|w|=R$ gives $|c_{2n}|\le M_RR^{-2n}$, so the resulting
series converges for every $y$ by taking $R^2>|y|$.
Its substitution and uniqueness follow coefficient by coefficient.
No square-root branch is introduced in this definition.

Set $f_\pm=F(\lambda_\pm)$ and $f'_\pm=F'(\lambda_\pm)$.
The exact first remainder and entire divided quotient are
\[
\begin{aligned}
 \ell_1&=(f_+-f_-)/d,&
 \ell_0&=(\lambda_+f_--\lambda_-f_+)/d,\\
 L(y)&=\ell_0+\ell_1y
 =\frac{(y-\lambda_-)f_+-(y-\lambda_+)f_-}{d},&
 Q(y)&=\frac{F(y)-L(y)}{H(y)},\\
 g(s)&=L((s-1/2)^2)+h(s)Q((s-1/2)^2).
\end{aligned}\tag{AUH3}
\]
The two singularities displayed in the quotient are removable because
$F-L$ vanishes at the two distinct nodes.  Division of its Taylor
series at each node proves removability and hence that $Q$ is entire.
It is jointly holomorphic in the distinct nodes as well: on any circle
enclosing the nodes and a local $y$ neighborhood its value is
\[
 Q(y)=\frac1{2\pi i}\int
 \frac{F(\zeta)}{(\zeta-\lambda_+)(\zeta-\lambda_-)(\zeta-y)}
 \,d\zeta.
 \tag{AUH4}
\]
The residues at the two nodes and $y$ give exactly AUH3;
coefficientwise integration proves joint holomorphic dependence.
Overlapping circles prove the assertion everywhere it is used.
Reality gives $\ell_0,\ell_1\in\mathbb R$ and
$Q(\bar y)=\overline{Q(y)}$ for the real original parameters.
In particular the original four values of $g$ vanish exactly when
$L=0$, equivalently $\ell_0=\ell_1=0$.
This equivalence follows because $L$ has degree at most one and
interpolates the two values; it does not presume that they vanish.

\subsection{The actual extended unit and all four first jets}

Define the two quotient values by their exact removable limits:
\[
 a_+=Q(\lambda_+)=\frac{f'_+-\ell_1}{d},\qquad
 a_-=Q(\lambda_-)=\frac{f'_--\ell_1}{-d},\qquad
 a_-=\overline{a_+}.
 \tag{AUH5}
\]
The last equality follows by conjugating $d$, $f'_+$ and the real
$\ell_1$.  These are actual values of the divided function AUH3 at
arbitrary original marks; they are not assigned amplitudes.
They are values of the original entire unit $v_h=g/h$ precisely on
the zero-residual locus $L=0$.  Away from it the exact identity is
$g/h=Q+L/H$, retaining the full meromorphic residual.

The four Xi jets specify every quantity in AUH3--5:
\[
\begin{aligned}
 f_+&=2\xi(\rho_{\mathrm M}),&f_-&=2\xi(\rho_{\mathrm N}),\\
 f'_+&=\xi'(\rho_{\mathrm M})/w_+,&
 f'_-&=\xi'(\rho_{\mathrm N})/w_-,\\
 a_+&=\frac{\xi'(\rho_{\mathrm M})/w_+
       -2(\xi(\rho_{\mathrm M})-\xi(\rho_{\mathrm N}))/d}{d},&
 a_-&=\frac{\xi'(\rho_{\mathrm N})/w_-
       -2(\xi(\rho_{\mathrm M})-\xi(\rho_{\mathrm N}))/d}{-d}.
\end{aligned}\tag{AUH6}
\]
The first line is AUH2 and the second is its derivative
$g'(1/2+w)=2wF'(w^2)$.  Neither $w_+$ nor $w_-$ is zero.
The map from the four scalar first-jet data to
$(\ell_0,\ell_1,a_+,a_-)$ is invertible, with exact inverse
\[
 f_\pm=\ell_0+\ell_1\lambda_\pm,\qquad
 f'_+=\ell_1+d a_+,\qquad f'_-=\ell_1-d a_-.
 \tag{AUH7}
\]
Its displayed inverse matrix has determinant $d^3\ne0$ in the
coordinate orders
\[
 (\ell_0,\ell_1,a_+,a_-),\qquad (f_+,f_-,f'_+,f'_-).
\]
Expanding its first two rows $(1,\lambda_\pm,0,0)$ and last two
rows $(0,1,d,0)$, $(0,1,0,-d)$ proves that determinant.
Thus the residual and quotient-unit coordinates together retain
all four original first-jet data.

At an actual simple quartet, AUH3 has $L=0$, so $Q((s-1/2)^2)=v_h(s)$
everywhere.  The unit diagonal and its value are then exactly
\[
 \bigl(v_h(\rho_{\mathrm M}),v_h(\rho_{\mathrm N}),
        v_h(\rho_{\mathrm T}),v_h(\rho_{\mathrm E})\bigr)
       =(a_+,a_-,a_-,a_+),\qquad
 a_+=\frac{2\xi'(\rho_{\mathrm M})}{h'(\rho_{\mathrm M})}
     =\frac{\xi'(\rho_{\mathrm M})}{4i\delta\gamma w_+}\ne0.
 \tag{AUH8}
\]
Indeed $h'(\rho_{\mathrm M})=2w_+d=8i\delta\gamma w_+$.
The nonvanishing follows exactly from simplicity; conversely, on
$L=0$ it makes all four zeros simple by reflection.
This proves that the original simple-quartet data are exactly
$\ell_0=\ell_1=0$, $a_+\ne0$ in the displayed real domain.
The entire function may have other zeros elsewhere; nothing in
this calculation asserts a global zero-free quotient.

\subsection{The full Hermite remainder modulo the original $h^2$}

Interpolate the two values in AUH5 a second time:
\[
 m_1=\frac{a_+-a_-}{d},\quad
 m_0=\frac{\lambda_+a_--\lambda_-a_+}{d},\quad
M(y)=m_0+m_1y,\quad
 R_2(y)=\frac{Q(y)-M(y)}{H(y)}.
 \tag{AUH9}
\]
Both $m_0,m_1$ are real and $R_2$ is entire by the same removable
division already proved.  The complete identity is
\[
\begin{aligned}
 g(s)&=J((s-1/2)^2)+h(s)^2R_2((s-1/2)^2),\\
 J(y)&=L(y)+H(y)M(y)=j_0+j_1y+j_2y^2+j_3y^3,\\
 (j_0,j_1,j_2,j_3)&=(\ell_0+\eta m_0,
       \ell_1+\beta m_0+\eta m_1,m_0+\beta m_1,m_1),\\
 [s^n]J((s-1/2)^2)&=
   \sum_{\substack{0\le r\le3\\2r\ge n}}
     \binom{2r}{n}(-1/2)^{2r-n}j_r\quad(0\le n\le7).
\end{aligned}\tag{AUH10}
\]
The last coefficient is zero when $n=7$, by the indicated empty sum.
Expansion proves every coefficient and preserves the original
$1/2$ translation with its signs.
The remainder is the unique polynomial of degree below eight in
$s$ matching $g$ and $g'$ at all four original marks.
It matches because $h^2$ has a double zero at every mark.
For uniqueness, the difference of two such polynomials is divisible
by each $(s-\rho_a)^2$, hence by their product $h^2$; its degree
is below eight, so it is zero.
Thus no part of the Hermite data has been discarded by using the
even presentation.  The identities AUH7 and
$a_\pm=m_0+m_1\lambda_\pm$ give the inverse from its coefficients
to the original jets; the determinant from
$(\ell_0,\ell_1,m_0,m_1)$ to those four jets is $-d^4\ne0$.

\subsection{The original four-dimensional primary multiplication matrix}

Work in the ordered centered coefficient basis $1,w,w^2,w^3$ of
$A_h=\mathbb C[w]/(w^4+\beta w^2+\eta)$.
Evaluation is the literal original matrix $V_{ar}=w_a^r$, with
rows $\mathrm M,\mathrm N,\mathrm T,\mathrm E$.
Its exact inverse columns and determinant are
\[
 (V^{-1})_{\cdot a}
 =\frac1{h_c'(w_a)}
 \begin{pmatrix}w_a^3+\beta w_a\\w_a^2+\beta\\w_a\\1\end{pmatrix},
 \quad h_c(w)=w^4+\beta w^2+\eta,\quad
 \det V=-64\delta^2\gamma^2(\delta^2+\gamma^2)\ne0.
 \tag{AUH11}
\]
To prove the inverse, the column is the coefficient vector of
$h_c(w)/((w-w_a)h_c'(w_a))$; division expands it exactly as
displayed.  It evaluates to one at $w_a$ and zero at the other
three marks.  This proves both inverse products.
The ordered Vandermonde product has factors
$-2i\gamma,-2\delta,-2w_+,-2w_-,-2\delta,-2i\gamma$,
giving the displayed determinant without reordering labels.
Its four derivative denominators in the same order are
$2w_+d,-2w_-d,2w_-d,-2w_+d$.

The original multiplier by the entire divided quotient $Q(w^2)$
in this algebra is exactly multiplication by $M(w^2)$:
\[
\begin{gathered}
 \mathsf M=
 \begin{pmatrix}
 m_0&0&-\eta m_1&0\\
 0&m_0&0&-\eta m_1\\
 m_1&0&m_0-\beta m_1&0\\
 0&m_1&0&m_0-\beta m_1
 \end{pmatrix},\qquad
 U_{\rm ext}=\operatorname{diag}(a_+,a_-,a_-,a_+),\\
 V\mathsf M V^{-1}=U_{\rm ext},\qquad
 \kappa=m_0^2-\beta m_0m_1+\eta m_1^2=a_+a_-=|a_+|^2,
 \qquad\det\mathsf M=\kappa^2.
\end{gathered}\tag{AUH12}
\]
Reducing $w^4=-\beta w^2-\eta$ and
$w^5=-\beta w^3-\eta w$ gives every matrix entry.
Evaluation commutes with multiplication, proving the similarity;
the product $(m_0+m_1\lambda_+)(m_0+m_1\lambda_-)$ proves
$\kappa$ and the determinant.
On the open set $a_+\ne0$, its exact inverse is obtained from
the same displayed matrix with
\[
 (m_0,m_1)\longmapsto
 (n_0,n_1)=\left(\frac{m_0-\beta m_1}{\kappa},
                         -\frac{m_1}{\kappa}\right),\qquad
 U_{\rm ext}^{-1}=\operatorname{diag}(a_+^{-1},a_-^{-1},a_-^{-1},a_+^{-1}).
 \tag{AUH13}
\]
In fact
$(m_0+m_1y)((m_0-\beta m_1)-m_1y)=\kappa-m_1^2H(y)$,
which proves the inverse in the quotient, and hence both matrix
inverse products.

For the original uncentered coefficient basis $1,s,s^2,s^3$, retain
$T_{rb}=\binom br2^{-(b-r)}$ for $r\le b$, zero otherwise.
It has determinant one, inverse
$T^{-1}_{rb}=\binom br(-1/2)^{b-r}$, and the matrices are exactly
\[
 V_\rho=VT,\qquad
 V_\rho^{-1}=T^{-1}V^{-1},\qquad
 \mathsf M_s=T^{-1}\mathsf M T,\qquad
 \mathsf M_s^{-1}=T^{-1}\mathsf M^{-1}T.
 \tag{AUH14}
\]
Substitution $s=1/2+w$ proves the two translation matrices compose
to the identity and establishes these formulas.
These are coordinate morphisms, not isometries.
For example if $H_{s_0,c}$ is the original Gamma Gram of the
uncentered monomials, the centered coefficient metric is
$(T^{-1})^*H_{s_0,c}T^{-1}$ and the primary evaluation metric is
$(V_\rho^{-1})^*H_{s_0,c}V_\rho^{-1}$.
Here $s_0\in\{1,a_{\rm amp}\}$ distinguishes the original
Gamma parameter from the polynomial variable.  Its mass, center
and all cross terms remain the original ones of WCF2 and RLA35;
no amplitude or metric factor has been set to one.

\subsection{The eight first-jet coordinates and the retained residual action}

Every polynomial class modulo $h_c^2$ has a unique expression
\[
 P_0(w)+h_c(w)P_1(w),\qquad \deg P_0,\deg P_1<4.
\]
Successive division by the monic quartic proves existence and
uniqueness.  Define $\mathsf L$ by AUH12's displayed matrix with
$(m_0,m_1)$ replaced by $(\ell_0,\ell_1)$, and let
$(\mathsf K_L)_{0,2}=(\mathsf K_L)_{1,3}=\ell_1$ be its only
nonzero entries.  The full multiplier by $g$ in these eight
coordinates is
\[
 \mathsf G_8=
 \begin{pmatrix}\mathsf L&0\\\mathsf M+\mathsf K_L&\mathsf L\end{pmatrix}.
 \tag{AUH15}
\]
Indeed multiplication of $L(w^2)$ by $P_0$ has reduced part
$\mathsf L P_0$ and the exact carry
$h_c\ell_1((P_0)_2+(P_0)_3w)$.
Multiplication of $h_cM(w^2)$ by $P_0$, and of $h_cL(w^2)$ by
$P_1$, gives the remaining lower block; every further carry has
a factor $h_c^2$ and is zero in the stated quotient.
This proves AUH15 from AUH10, including the residual carry which
would be lost by replacing the nonzero residual by a unit.

Let $V'_{ar}=r w_a^{r-1}$, with its column at $r=0$ equal to zero,
and let $D_h=\operatorname{diag}(h_c'(w_a))$ in the unchanged order.
The exact evaluation-and-derivative map, and its inverse, are
\[
\begin{aligned}
 \mathsf J_8&=\begin{pmatrix}V&0\\V'&D_hV\end{pmatrix},&
 \binom{\mathbf v}{\mathbf j}&=\mathsf J_8\binom{P_0}{P_1},\\
 P_0&=V^{-1}\mathbf v,&
 P_1&=V^{-1}D_h^{-1}(\mathbf j-V'V^{-1}\mathbf v),\\
 \mathsf J_8\mathsf G_8\mathsf J_8^{-1}
 &=\begin{pmatrix}D_g&0\\D_{g'}&D_g\end{pmatrix},&
 (D_g)_{aa}&=L(w_a^2),\quad
 (D_{g'})_{aa}=2w_a\ell_1+h_c'(w_a)\,a_a,\\
 (a_{\mathrm M},a_{\mathrm N},a_{\mathrm T},a_{\mathrm E})
 &=(a_+,a_-,a_-,a_+).
\end{aligned}\tag{AUH16}
\]
The first line is the derivative of $P_0+h_cP_1$ at each root;
substitution verifies both inverse compositions.  The third line
is the product rule $(gP)'=g'P+gP'$, with $g,g'$ from AUH10.
It proves the full eight-jet identity, not just its four-value
restriction.  If both sampled values $f_+,f_-$ are nonzero, then
$\mathsf L$ is inverted by AUH13 with the $\ell$ coefficients,
and the entire eight-dimensional inverse is
\[
 \mathsf G_8^{-1}=
 \begin{pmatrix}
 \mathsf L^{-1}&0\\
 -\mathsf L^{-1}(\mathsf M+\mathsf K_L)\mathsf L^{-1}&\mathsf L^{-1}
 \end{pmatrix}.
 \tag{AUH17}
\]
At the original simple quartet, instead, $L=0$ and AUH15 is
$\left(\begin{smallmatrix}0&0\\\mathsf M&0\end{smallmatrix}\right)$,
with square zero and rank four, since $\mathsf M$ is invertible.
Thus the unit and the nilpotent multiplication by $g$ have their
different exact maps retained.  Formula AUH16 remains valid at
every stratum, whether or not AUH17 is invertible.

\subsection{Two exact phase charts covering every nonzero actual unit}

Put $r=(a_++a_-)/2=\Re a_+$ and
$c=(a_+-a_-)/(2i)=\Im a_+$.
They are determined by AUH6, not free coordinates in the pullback.
Write $J_d=\operatorname{diag}(1,-1,-1,1)$, so
$U_{\rm ext}=rI+icJ_d$.
The complete two-chart dictionary and its inverses are
\[
\begin{array}{c|c|c|c}
 \text{chart}&\text{coordinates}&U_{\rm ext}&\text{matrix inverse}\\\hline
 r\ne0&(r,x=c/r)&rW_R(x),\quad W_R=I+ixJ_d&
 W_R^{-1}=(I-ixJ_d)/(1+x^2)\\
 c\ne0&(c,y=r/c)&cW_I(y),\quad W_I=yI+iJ_d&
 W_I^{-1}=(yI-iJ_d)/(1+y^2)
\end{array}.
 \tag{AUH18}
\]
The inverse coordinate formulas are $a_+=r(1+ix)$ and
$a_+=c(y+i)$, with $a_-$ its conjugate.
On the overlap,
\[
 y=1/x,\qquad c=rx,\qquad r=cy,\qquad
 W_I(1/x)=x^{-1}W_R(x).
 \tag{AUH19}
\]
Direct multiplication proves the matrix inverses and transitions.
Since $a_+\ne0$ exactly when $(r,c)\ne(0,0)$, these two charts
cover every nonzero actual unit, including a purely imaginary one.
The original unit magnitude is $r^2+c^2=|a_+|^2=\kappa$,
not a chosen unit modulus.  It cancels only where an operator
contains $U_{\rm ext}^{-1}$ and $U_{\rm ext}$ as a conjugate pair.
All period matrices and original metrics still retain the full
$U_{\rm ext}$.

\subsection{The exact arithmetic-jet morphism into the original conductor family}

For the original nonzero period $u$ on its unchanged logarithm branch,
put $z=u^{-1}$, $\mathfrak a=1/160+\beta/24+\eta/2$,
$\zeta=e^{2\pi i/5}$ and
$D_j=\operatorname{diag}(\zeta^{-ja})_{a=1}^4$.
Use the entire original period matrix, with every coefficient,
\[
\begin{aligned}
 \mathcal R_{ar}(z;\delta,\gamma)
 &=\sum_{n\ge0}\ 
 \sum_{\substack{p,h\ge0\\2p+4h=5n+r+1-a}}
 \frac{(\beta/3)^p\eta^h(-5)^{p+h-n}(a/5)_{p+h-n}}{p!h!}\,z^n,\\
 \mathsf H_{\rm ext}(u)&=
 e^{\mathfrak a/u}G_u\mathcal R(u^{-1};\delta,\gamma)V^{-1}U_{\rm ext},\\
 G_u&=\operatorname{diag}(u^{a/5}\,\mathfrak g_a)_{a=1}^4,\qquad
 \mathfrak g_a=5^{a/5-1}e^{\pi ia/5}\Gamma(a/5),\\
 \mathsf C_j&=V\mathcal R^{-1}D_j\mathcal RV^{-1},\qquad
 \mathsf B_j=U_{\rm ext}^{-1}\mathsf C_jU_{\rm ext}.
\end{aligned}\tag{AUH20}
\]
The index ranges are $1\le a\le4$, $0\le r\le3$, and empty
inner sums are zero.  PCL1--5 prove convergence and
$\det\mathcal R=1$, with its exact inverse $\operatorname{adj}\mathcal R$.
PQO8--9 and PCL3 specify the retained Gamma ray factors $\mathfrak g_a$;
every power of $u$ in $G_u$ uses the same original branch.
The original conjugation identity is exactly
\[
 \mathsf H_{\rm ext}^{-1}D_j\mathsf H_{\rm ext}=\mathsf B_j.
 \tag{AUH21}
\]
Indeed the nonzero scalar exponential cancels with its inverse and
$G_u$ commutes with $D_j$.  The inverse of the entire left-hand
coordinate map, where $a_+\ne0$, is explicitly
$U_{\rm ext}^{-1}V\mathcal R^{-1}G_u^{-1}e^{-\mathfrak a/u}$.
No Gamma factor or center is removed from that map.

Define the source parameter domain
$\mathcal D=\{(\delta,\gamma,u):0<\delta<1/2,\ \gamma>2,
\ u\ne0,\ a_+(\delta,\gamma)\ne0\}$,
where $a_+$ is the actual Xi-jet expression AUH6.
The following are exact morphisms into the extended original
geometric family, with the residual carried as additional output:
\[
\begin{aligned}
 \Psi:(\delta,\gamma,u)&\longmapsto
  (\delta,\gamma,z=u^{-1},a_+,a_-;\ell_0,\ell_1),\\
 \Psi_R:(\delta,\gamma,u)&\longmapsto
  (\delta,\gamma,z=u^{-1},r,x=c/r;\ell_0,\ell_1)
                                      &&(r\ne0),\\
 \Psi_I:(\delta,\gamma,u)&\longmapsto
  (\delta,\gamma,z=u^{-1},c,y=r/c;\ell_0,\ell_1)
                                      &&(c\ne0).
\end{aligned}\tag{AUH22}
\]
Their inverse onto their respective graphs retains $\delta,\gamma$
and takes $u=1/z$; AUH18--19 recover every unit and residual
coordinate.  Thus these are not freely assigned phase families.
The first chart uses exactly RLA2's $W(x)$; the second gives its
regular complementary chart by the proved transition.
The quartet parameters in AUH20 are retained as variables; fixing
them at $(1/4,3)$ gives precisely RLA2's coefficient family.
The isomorphism of local branches proved at the RLA base point is
not asserted at every image of AUH22.  AUH22 itself is an exact
map on its entire stated domain, requiring no assumed zeros.

The actual simple-quartet locus in this domain is exactly
$\mathcal Z=\{\ell_0=\ell_1=0\}$, by AUH3 and AUH8.
On its complement AUH22 still composes the original geometric
matrix formula with the actual divided Xi-jet values, and retains
the exact discrepancy $L/h$ from the original quotient.
Consequently the map does not label that nonzero-residual extension
as an arithmetic simple-quartet unit.  This distinction is given
by the explicit equality AUH3, not by an unconstructed comparison.

\subsection{Every pulled-back orbit factor and the full zeroth moment}

To distinguish the original Riemann function from the exponential
spectral coordinate, write that latter coordinate as $\tau$ in this
subsection; it is exactly the variable called $\xi$ in RLA2.
With $e(\tau)=(e^{\rho_a\tau})_a$, the complete composition is
\[
\begin{aligned}
 f_j^{\Xi}(\tau;\delta,\gamma,u)
 &=Q_0\bigl(\mathsf B_j e(\tau)\bigr),\qquad Q_0(v)=v_1v_4-v_2v_3,\\
 E_A^{\Xi}(\tau;\delta,\gamma,u)&=\prod_{j=1}^4f_j^{\Xi}(\tau),\qquad
 \mu_n^{\Xi}=\partial_\tau^nE_A^{\Xi}(0),\\
 g_{\rm cond}^{\Xi}(t)&=e^{-4\omega(t)}E_A^{\Xi}(\omega(t)),
       \qquad\omega(t)=-i\arctan t.
\end{aligned}\tag{AUH23}
\]
Every entry of $\mathsf B_j$ is the convergent matrix expression
AUH20 with the actual Xi samples AUH6.  It is therefore an explicit
composition, with every original ordered-pair multiplicity retained
by the products.  The original symbol in WCF and RLA, its $-4$
center correction, every moment and the unchanged Gamma weights
are all pulled back by this identity.

There is also an exact finite expression displaying the amplitude
dependence without concealing the residual constraint.
Let $e_+=(e^{\rho_{\mathrm M}\tau},0,0,e^{\rho_{\mathrm E}\tau})^{\mathsf T}$,
$e_-=(0,e^{\rho_{\mathrm N}\tau},e^{\rho_{\mathrm T}\tau},0)^{\mathsf T}$,
and write $p_{jr}=(\mathsf C_je_+)_r$,
$q_{jr}=(\mathsf C_je_-)_r$.
For $\vartheta=a_-/a_+$, the full factor is
\[
\begin{aligned}
 f_j^{\Xi}
 ={}&p_{j1}p_{j4}
 +\vartheta(p_{j1}q_{j4}+q_{j1}p_{j4})
 +\vartheta^2q_{j1}q_{j4}\\
 &-\vartheta^{-2}p_{j2}p_{j3}
 -\vartheta^{-1}(p_{j2}q_{j3}+q_{j2}p_{j3})
 -q_{j2}q_{j3}.
\end{aligned}\tag{AUH24}
\]
This follows by applying $U_{\rm ext}$ to $e_++e_-$ and
$U_{\rm ext}^{-1}$ to the four output coordinates in AUH20;
in particular its four signs and its two cross terms are exact.
The common amplitude scale cancels here because it occurs in a
conjugate matrix pair; its value remains in AUH20's coordinate map.
It has not been assigned by selecting $\vartheta$.

An expression without any phase chart is
\[
\begin{aligned}
 \mathcal D_j(\tau)
 ={}&a_-^2(a_+p_{j1}+a_-q_{j1})(a_+p_{j4}+a_-q_{j4})\\
 &-a_+^2(a_+p_{j2}+a_-q_{j2})(a_+p_{j3}+a_-q_{j3}),\\
 f_j^{\Xi}(\tau)&=\frac{\mathcal D_j(\tau)}{a_+^2a_-^2},\qquad
 \mu_0^{\Xi}(\delta,\gamma,u)
 =\frac{\prod_{j=1}^4\mathcal D_j(0)}{(a_+a_-)^8}.
\end{aligned}\tag{AUH25}
\]
Both equalities follow by clearing exactly the displayed denominators
in AUH24.  The numerators remain defined even at vanishing quotient
values; the inverse coordinate map at such a value is not asserted.
For real $u\ne0$, the original reflection calculation RLA2 gives
$f_{5-j}^{\Xi}(\tau)=-\overline{f_j^{\Xi}(\bar\tau)}$.
Indeed the coefficient series is real, $\bar V=J_cV$,
$\bar U_{\rm ext}=J_cU_{\rm ext}J_c$ and
$Q_0(J_cv)=-Q_0(v)$, exactly as in the original proof.
Therefore the exact real-parameter pullback is
\[
 \mu_0^{\Xi}
 =|f_1^{\Xi}(0)f_2^{\Xi}(0)|^2
 =\frac{|\mathcal D_1(0)\mathcal D_2(0)|^2}{(a_+a_-)^8}.
 \tag{AUH26}
\]
This is the complete zeroth moment on the actual Xi-jet graph.
Its zero set there is exactly the union of the two explicit
numerator zero sets.  The arithmetic simple-quartet intersection
also retains the two explicit equations $\ell_0=\ell_1=0$.
These are pullbacks of the actual functions, rather than a claim
that an arbitrary geometric phase pair meets that intersection.

\subsection{Exact differentials of the pullback, including all Xi and period terms}

To make the parameter morphism usable in local calculations, retain
ordinary differentials in the original variables:
\[
\begin{aligned}
 dw_\pm&=d\delta\pm i\,d\gamma,\qquad
 d\lambda_\pm=2w_\pm dw_\pm,\\
 d(d)&=4i(\gamma\,d\delta+\delta\,d\gamma),\\
 d\ell_1&=\frac{f'_+d\lambda_+-f'_-d\lambda_--\ell_1d(d)}{d},\\
 d\ell_0&=df_+-\lambda_+d\ell_1-\ell_1d\lambda_+,\\
 da_+&=\frac{F''(\lambda_+)d\lambda_+-d\ell_1-a_+d(d)}{d},\\
 da_-&=\frac{F''(\lambda_-)d\lambda_--d\ell_1+a_-d(d)}{-d},\\
 F''(\lambda_\pm)&=
   \frac{\xi''(\rho_{\mathrm M/\mathrm N})}{2w_\pm^2}
  -\frac{\xi'(\rho_{\mathrm M/\mathrm N})}{2w_\pm^3},\\
 dz&=-u^{-2}du.
\end{aligned}\tag{AUH27}
\]
Here $d(d)$ explicitly means the differential of the scalar
$d=4i\delta\gamma$; $df_+=f'_+d\lambda_+$.
Each identity is differentiation of AUH3--6; the second derivative
formula follows by differentiating $g'(1/2+w)=2wF'(w^2)$.
Thus every second Xi derivative needed for the actual phase
pullback is specified.  The exact chart derivatives are
\[
 dr=(da_++da_-)/2,\quad dc=(da_+-da_-)/(2i),\quad
 dx=(r\,dc-c\,dr)/r^2,\quad
 dy=(c\,dr-r\,dc)/c^2.
 \tag{AUH28}
\]
They apply on the respective charts of AUH18, and retain the
actual scale derivative rather than holding a free phase fixed.

Finally let $K=\mathcal RV^{-1}U_{\rm ext}$, so that
$\mathsf B_j=K^{-1}D_jK$.
Every parameter derivative of the pulled-back original matrices is
\[
\begin{aligned}
 \mathcal A=K^{-1}dK
 &=U_{\rm ext}^{-1}V\mathcal R^{-1}(d\mathcal R)V^{-1}U_{\rm ext}
   -U_{\rm ext}^{-1}(dV)V^{-1}U_{\rm ext}
   +U_{\rm ext}^{-1}dU_{\rm ext},\\
 d\mathsf B_j&=[\mathsf B_j,\mathcal A],\qquad
 (dV)_{ar}=r w_a^{r-1}dw_a,\\
 d\mathcal R&=\mathcal R_zdz+\mathcal R_\beta d\beta+
                                    \mathcal R_\eta d\eta,\\
 d\beta&=4\gamma\,d\gamma-4\delta\,d\delta,\qquad
 d\eta=4(\delta^2+\gamma^2)(\delta\,d\delta+\gamma\,d\gamma),\\
 dU_{\rm ext}&=\operatorname{diag}(da_+,da_-,da_-,da_+).
\end{aligned}\tag{AUH29}
\]
The coefficients of $\mathcal R_\beta,\mathcal R_\eta$ are the
literal term derivatives of AUH20: a term indexed by $(p,h)$ is
multiplied respectively by $p/\beta$ and $h/\eta$, with zero
for $p=0$ or $h=0$.
Both $\beta,\eta$ are positive on the real original domain.
Termwise differentiation is justified on compact parameter sets by
the original factorially convergent period series, or by its
absolutely convergent period integral with the fifth-degree decay
and polynomial parameter derivatives.  For an explicit majorant,
choose a common constant $K_0\ge1$ on the compact set with
$K_0^2\ge5|\beta|/3$, $K_0^4\ge5|\eta|$.
The coefficient argument of RPZ5 \cite{RPZ005} bounds the entire coefficient
by $4K_0^3(2n+2)(6K_0^5/5)^n/n!$.
On a slightly larger compact set with nonzero $\beta,\eta$ the
first coefficient derivatives introduce factors at most linear in
$n$, because $p,h\le(5n+3)/2$; higher fixed derivatives introduce
fixed polynomial factors.  The resulting series still converge
uniformly.  This supplies the differentiation used in AUH29.
Differentiation of $K^{-1}D_jK$ proves the commutator, including
its order and all three terms in $\mathcal A$.

Let $v_j=\mathsf B_je(\tau)$ and retain the original bilinear
polar form
$\mathfrak b(v,w)=v_1w_4+v_4w_1-v_2w_3-v_3w_2$.
Then the full first variation is
\[
\begin{aligned}
 df_j^{\Xi}&=\mathfrak b\bigl(v_j,
        [\mathsf B_j,\mathcal A]e(\tau)+\mathsf B_jde(\tau)\bigr),\\
 (de(\tau))_a&=e^{\rho_a\tau}(\rho_a\,d\tau+\tau\,d\rho_a),\\
 d\mu_0^{\Xi}&=\sum_{j=1}^4df_j^{\Xi}(0)
                               \prod_{k\ne j}f_k^{\Xi}(0).
\end{aligned}\tag{AUH30}
\]
The first line follows from $Q_0(v)=\mathfrak b(v,v)/2$ and
symmetry of this bilinear form, with no Hermitian substitution.
The last line is the exact product rule, valid also at its zero
set without division by a vanishing factor.
Equations AUH27--30 give the actual tangent pullback into RLA's
factor and moment equations, including changes of the quartet,
unit, period and original spectral coordinate.

This completes the morphism from the actual Xi-jet data to the
original conductor and its local algebra wherever the latter's
coordinate chart is used.  The full residual, the unit scale and
both phase charts remain explicit.  The calculation proves the
map and its identities; it does not assert that the previously
certified geometric singular point is on this arithmetic graph.
